\begin{adjustwidth}{-\extralength}{0cm}
\begin{longtblr}[
  label={tab:criterios_ie},
  caption={Inclusion and exclusion criteria applied in the \textit{screening} and eligibility phases, aligned with PICo, research directives (D1--D3), and Kitchenham's practical criteria. Own elaboration.}]
{
    colspec = {X[m,j,3] X[m,j,5] X[m,j,5] X[m,j,1.25]},
    width=\fulllength,
    rowhead=1,
    row{1}={font=\bfseries}
}
\toprule
\textbf{Criterion} &
\textbf{Inclusion (Operational description)} &
\textbf{Exclusion (Operational description)} &
\textbf{Phase} \\
\midrule

Relation to the phenomenon (Air quality -- P) &
The study explicitly addresses air quality or ambient atmospheric pollution. &
The study focuses on water, soil, agriculture, aquaculture, noise, or another domain unrelated to air quality. &
Screening \\

Central use of fuzzy logic (I) &
The title or abstract mentions \textit{fuzzy logic} or fuzzy inference systems as the main approach. &
No use of fuzzy logic is identified, or it does not constitute the study's central approach. &
Screening \\

Air-quality evaluation through fuzzy logic (P + I) &
The study proposes evaluation, classification, or analysis of air-quality level using fuzzy logic. &
Fuzzy logic is applied only to control (PID, HVAC, drying, ventilation, furnaces, etc.) without evaluating air quality as the main variable. &
Screening \\

Nature of the study (Article type) &
Primary study with an original technical proposal or implementation. &
Systematic review, narrative review, \textit{mapping study}, \textit{proceedings} volume, or grey literature. &
Screening \\

Air-quality monitoring platform or system (Co) &
The study proposes or implements a platform, monitoring system, IoT architecture, sensor integration, or structured use of \textit{datasets} for air-quality analysis. &
The study does not describe a system, platform, or data integration, or is limited to theoretical analysis without a monitoring or structured processing component. &
Screening \\

\textbf{(D1) }Specification of variables and pollutants used &
The study explicitly identifies the atmospheric pollutants and/or environmental variables used as inputs (e.g., PM$_{2.5}$, PM$_{10}$, NO$_2$, O$_3$, CO, SO$_2$, temperature, humidity, wind). &
The variables used are not clearly specified, or they cannot be identified from the full text. &
Full Text \\

\textbf{(D2) }Explicit application of a fuzzy inference system to air quality &
Fuzzy inference constitutes the central mechanism for evaluating air quality and/or associated risk. &
Fuzzy logic is tangential, secondary, or not directly applied to air-quality evaluation. &
Full Text \\

\textbf{(D2) }Structural description of the fuzzy model (type and design) &
The system type (e.g., Mamdani, Sugeno, ANFIS, type-2, or other) is specified, and structural elements (rules, membership functions, or architecture) are described. &
Neither the fuzzy-system type nor its structure is detailed, preventing characterization and systematic extraction. &
Full Text \\

\textbf{(D3) }Implementation in a monitoring system or platform &
The study describes implementation in a platform, IoT architecture, monitoring system, or structured air-quality data-processing environment. &
The model is presented without the context of a monitoring system/platform or without operational integration into a data-processing environment. &
Full Text \\

\textbf{(D3) }Description of application setting and unit of analysis (Setting / Participants) &
The application context (e.g., urban/city, stations, sensors, \textit{datasets}) and the evaluated unit of analysis are specified. &
The application setting or unit of analysis is not clearly described, preventing comparison with other studies. &
Full Text \\

\textbf{(Cross-cutting) }Documented methodological design and validation (Research Design / Sampling) &
The study describes an experimental or applied design and presents empirical validation or quantitative evaluation (e.g., metrics, comparison, test cases). &
It does not present a clear methodological design and/or does not report empirical validation or quantitative evaluation of the proposed system. &
Full Text \\

\bottomrule
\end{longtblr}
\end{adjustwidth}
